\chapter{Architecture and planning}
\label{ch::chapter3}

\section{Introduction}
This chapter presents the technical foundation of Printymand. It defines the actors, functional and non-functional requirements, global architecture, database model, front-end and back-end organization, and the work environment used during development. The chapter also maps the project work into Spiral cycles.

\section{Requirements specification}
\subsection{Identification of actors}
Printymand is a multi-role platform. The main actors are:

\begin{itemize}
    \item \textbf{Visitor:} explores public pages and can access authentication pages.
    \item \textbf{Customer:} browses the marketplace, customizes products, selects printers, manages cart content, places orders, tracks fulfillment, and submits reviews.
    \item \textbf{Designer:} manages profile information, uploads or edits designs, assigns products to designs, and views sales and engagement indicators.
    \item \textbf{Printer:} manages printable products, receives order lines, and updates production statuses such as pending, accepted, printing, shipped, delivered, or rejected.
    \item \textbf{Administrator:} supervises users, rankings, content, dashboard indicators, and system health.
\end{itemize}

\subsection{Functional requirements}
The main functional requirements deduced from the codebase are:

\begin{itemize}
    \item Allow users to register and log in according to a role: customer, designer, printer, or administrator.
    \item Protect private routes through authentication and role-based authorization.
    \item Display a home page, marketplace page, product catalog, and design detail page.
    \item Allow customers to customize a selected design on an available product by changing color, size, position, and scale.
    \item Allow customers to select a printer before adding the customized item to the cart.
    \item Allow cart review, checkout, payment-mode selection, and local order creation.
    \item Allow order tracking with a timeline based on order-line statuses.
    \item Provide customer profile management, saved addresses, payment preferences, and order history.
    \item Provide designer dashboards for design management, product assignment, analytics, and payout overview.
    \item Provide printer dashboards for product management and fulfillment status updates.
    \item Provide administrator dashboards for global indicators, users, roles, ranks, and content supervision.
    \item Provide back-end endpoints for authentication, designs, products, dashboard data, health checks, and future cart operations.
\end{itemize}

\subsection{Non-functional requirements}
The following non-functional requirements are relevant to Printymand:

\begin{itemize}
    \item \textbf{Security:} passwords are hashed with bcrypt; API authentication uses JWT; protected endpoints use guards; authentication tokens are stored in HTTP-only cookies.
    \item \textbf{Maintainability:} the back end follows a modular layered structure with controllers, use cases, repositories, domain DTOs, and infrastructure services.
    \item \textbf{Usability:} the front end provides role-specific pages, clear navigation, dashboard shells, reusable cards, and forms.
    \item \textbf{Scalability:} PostgreSQL schema separates users, inventories, products, designs, tags, orders, and order items to support marketplace growth.
    \item \textbf{Reliability:} the back end includes global exception filters, validation pipes, logging interceptors, transactions, and health checks.
    \item \textbf{Performance:} cursor pagination structures are prepared for designs and products; throttling is configured to limit abusive request rates.
\end{itemize}

\section{Global analysis}
\subsection{Global use case diagram}
The global use case diagram should represent all main interactions between customers, designers, printers, administrators, and the Printymand system. Since the real diagram is not present in the repository, the report includes a placeholder path for later replacement.

\placeholderfigure{global-use-case}{figures/global-use-case-diagram.png}{Global use case diagram placeholder.}

\subsection{Global class diagram}
The main entities identified from the codebase are User, CustomerProfile, DesignerProfile, PrinterProfile, Design, Product, PrinterPartner, CustomizationDraft, CartItem, Order, OrderLine, Review, and PayoutRecord. On the back end, the PostgreSQL migrations define users, designer inventories, designer albums, designs, design tags, printer inventories, products, orders, and order items.

\placeholderfigure{global-class-diagram}{figures/global-class-diagram.png}{Global class diagram placeholder.}

\subsection{Global sequence diagram}
The global sequence starts when a customer authenticates, browses active designs, selects a product, customizes placement, chooses a printer, adds the draft to the cart, places an order, and follows the tracking timeline. Designers and printers interact with the system through role dashboards, while the administrator supervises the platform.

\placeholderfigure{global-sequence-diagram}{figures/global-sequence-diagram.png}{Printymand global sequence diagram placeholder.}

\section{System architecture}
Printymand follows a client-server architecture. The Angular application is the user-facing single-page application. It communicates with a NestJS API through HTTP requests under the \texttt{/api/v1} prefix. The back end persists data in PostgreSQL and exposes structured modules for authentication, products, designs, dashboard, health, users, and cart.

\placeholderfigure{system-architecture}{figures/system-architecture-diagram.png}{System architecture diagram placeholder.}

\subsection{Front-end architecture}
The front end is located in the \texttt{frontend} folder and is implemented with Angular 21 and TypeScript. It uses standalone components, Angular routing, signals, computed state, services, and template files. The routing file defines pages for:

\begin{itemize}
    \item Home, marketplace, products, and design details.
    \item Login and registration.
    \item Customization, printer selection, checkout, and order tracking.
    \item Customer, designer, printer, and administrator dashboards and profiles.
    \item A not-found page for invalid routes.
\end{itemize}

The front end contains reusable components such as layout, dashboard shell, design card, product card, printer card, stat card, rank badge, image-with-fallback, skeleton block, and UI modal. This organization encourages reuse and separates page-level workflows from common interface elements.

The \texttt{PlatformStoreService} is central to the current front-end implementation. It stores users, printers, products, designs, cart items, orders, reviews, payouts, and customization drafts. It also persists state in browser local storage. The service currently supports three modes: mock, hybrid, and API. This is important because the front-end experience is already broad, while parts of the back-end marketplace API are still being completed.

\subsection{Back-end architecture}
The back end is located in the \texttt{backend} folder and is implemented with NestJS 11 and TypeScript. Its structure follows a layered approach:

\begin{itemize}
    \item \textbf{API layer:} controllers, guards, decorators, filters, interceptors, and query DTOs.
    \item \textbf{Application layer:} use cases such as login, register, logout, refresh, view products, get product, add product, view designs, get design, upload design, and get dashboard data.
    \item \textbf{Domain layer:} entities, DTOs, repository contracts, policies, and value objects.
    \item \textbf{Infrastructure layer:} JWT service, persistence module, PostgreSQL service, repository implementations, transactions, realtime services, and health checks.
\end{itemize}

The back end configures CORS for the Angular development server, uses cookie parsing, URI versioning, a global \texttt{/api} prefix, cache manager, throttling, validation pipes, class serialization, logging interceptors, and global exception filters.

\subsection{Authentication and authorization}
Authentication is implemented through the \texttt{AuthController}. It exposes registration, login, logout, current-user, and refresh endpoints. Passwords are hashed and verified with bcrypt through a password policy. JWT tokens are signed by the back-end JWT service and written to an HTTP-only cookie. The \texttt{JwtGuard} extracts tokens from cookies or authorization headers. The \texttt{RbacGuard} checks role metadata declared with the \texttt{Roles} decorator.

The front end also protects routes through an \texttt{authRoleGuard}. Role-specific pages such as designer dashboard, printer dashboard, and admin dashboard are guarded according to the current user role.

\subsection{Database usage}
The PostgreSQL migrations define the initial relational model:

\begin{itemize}
    \item \texttt{users}: stores account identity, password hash, contact data, role, and JSON data.
    \item \texttt{printer\_inventory} and \texttt{product}: represent printer-owned printable items.
    \item \texttt{designer\_inventory}, \texttt{designer\_album}, \texttt{design}, \texttt{design\_tag}, and \texttt{design\_tags}: represent designer assets and categorization.
    \item \texttt{order} and \texttt{order\_item}: represent placed orders and their line items.
\end{itemize}

The persistence layer uses the \texttt{pg} package through a \texttt{PostgresService}. A transaction scope is provided to execute operations inside database transactions. Repository implementations map SQL query results into DTOs.

\subsection{API routes and modules}
The current API modules include:

\begin{itemize}
    \item \texttt{/auth/register}, \texttt{/auth/login}, \texttt{/auth/logout}, \texttt{/auth/me}, and \texttt{/auth/refresh}.
    \item \texttt{/products} for product creation by printers and product listing/detail retrieval.
    \item \texttt{/designs} for designer design actions and design listing/detail retrieval.
    \item \texttt{/dashboard} for dashboard data.
    \item \texttt{/health} for administrator-only health checks.
    \item \texttt{/cart} as a module prepared for cart-related operations.
\end{itemize}

Some endpoints and repositories are intentionally still incomplete in the current development state. For example, the cart controller is a placeholder, product SQL queries are not fully implemented, and front-end API paths still include some legacy route expectations. This confirms that the project is an academic in-development platform rather than a finished production system.

\section{Implementation roadmap}
\begin{longtable}{|p{0.12\linewidth}|p{0.25\linewidth}|p{0.45\linewidth}|p{0.12\linewidth}|}
\caption{Implementation roadmap.} \label{tab:implementation_roadmap} \\
\hline
\textbf{Cycle} & \textbf{Focus} & \textbf{Main deliverables} & \textbf{Status} \\ \hline
\endfirsthead
\hline
\textbf{Cycle} & \textbf{Focus} & \textbf{Main deliverables} & \textbf{Status} \\ \hline
\endhead
\hline
\endfoot
1 & Authentication and discovery & Registration, login, role guards, marketplace browsing, product and design pages, reusable UI components. & Implemented mainly in front end, back-end authentication started. \\ \hline
2 & Customization and checkout & Customization draft, product preview, printer selection, cart, order placement, payment-mode flow, tracking. & Implemented in front-end mock/hybrid workflow. \\ \hline
3 & Governance and optimization & Customer, designer, printer, and admin dashboards; PostgreSQL schema; health checks; validation; RBAC; backend modularization. & Partially implemented and under consolidation. \\ \hline
\end{longtable}

\section{Work environment}
\subsection{Software architectural pattern}
The front end follows a component-service architecture typical of Angular applications. Pages handle view-level concerns, reusable components encapsulate interface elements, and services centralize authentication, API communication, and workflow state.

The back end follows a clean layered architecture. Controllers receive HTTP requests, use cases implement application logic, repository interfaces define persistence contracts, and infrastructure classes implement technical details such as PostgreSQL access and JWT generation.

\subsection{Modeling language}
The Unified Modeling Language (UML) is used to represent use cases, class diagrams, and sequence diagrams. UML is appropriate because it helps describe actors, entities, relationships, and workflows in a way that is understandable for both technical and academic audiences.

\subsection{Hardware environment}
\begin{table}[!htpb]
\caption{Hardware environment table.}
\centering
\begin{tabular}{|p{0.35\linewidth}|p{0.5\linewidth}|}
\hline
\textbf{Element} & \textbf{Specification} \\ \hline
Development machines & Student laptops used for front-end and back-end development \\ \hline
Operating system & Windows development environment \\ \hline
Memory & Suitable for Angular, NestJS, PostgreSQL, and LaTeX authoring \\ \hline
Network & Local development environment using Angular and NestJS servers \\ \hline
\end{tabular}
\label{tab:hardware_env}
\end{table}

\subsection{Software environment}
\begin{longtable}{|p{0.25\linewidth}|p{0.62\linewidth}|}
\caption{Development tools table.} \label{tab:development_tools} \\
\hline
\textbf{Tool} & \textbf{Usage in the project} \\ \hline
\endfirsthead
\hline
\textbf{Tool} & \textbf{Usage in the project} \\ \hline
\endhead
\hline
\endfoot
Visual Studio Code & Source code editing and project navigation. \\ \hline
Git and GitHub & Version control and collaboration. \\ \hline
Node.js and npm & Package management and execution of Angular and NestJS scripts. \\ \hline
PostgreSQL & Relational database for the back-end persistence layer. \\ \hline
LaTeX & Professional report writing and academic document generation. \\ \hline
\end{longtable}

\begin{longtable}{|p{0.28\linewidth}|p{0.58\linewidth}|}
\caption{Libraries, frameworks and programming languages table.} \label{tab:libraries_frameworks} \\
\hline
\textbf{Technology} & \textbf{Usage in the project} \\ \hline
\endfirsthead
\hline
\textbf{Technology} & \textbf{Usage in the project} \\ \hline
\endhead
\hline
\endfoot
Angular 21 & Single-page front-end application, routing, components, forms, signals, and services \cite{Angular}. \\ \hline
TypeScript & Main language for front-end and back-end development \cite{TypeScript}. \\ \hline
NestJS 11 & Back-end framework used for modular API development, controllers, guards, providers, and dependency injection \cite{NestJS}. \\ \hline
PostgreSQL & Relational database used through SQL migrations and the \texttt{pg} driver \cite{PostgreSQL}. \\ \hline
JWT & Authentication token format used by the back end \cite{JWT}. \\ \hline
bcrypt & Password hashing and verification library. \\ \hline
Socket.IO & Dependency prepared for realtime services and presence features. \\ \hline
Tailwind CSS & Front-end styling dependency used in the Angular project \cite{Tailwind}. \\ \hline
\end{longtable}

\section{Conclusion}
This chapter described the requirements, architecture, actors, database usage, API modules, and development environment of Printymand. It also clarified the current maturity of the platform: the front end already demonstrates most workflows, while the back end provides the foundations and progressively implements marketplace operations. The next chapter presents the first Spiral cycle.
